% 陈成：其他个人经历与特长
% 使用 XeLaTeX 编译：xelatex cheng-chen-additional-profile-zh.tex
\documentclass[UTF8,fontset=none,10pt,a4paper]{ctexart}

\usepackage[left=1.9cm,right=1.9cm,top=1.55cm,bottom=1.55cm]{geometry}
\usepackage{enumitem}
\usepackage{fontspec}
\usepackage{hyperref}
\usepackage{tabularx}
\usepackage{xcolor}

\setmainfont{Times New Roman}
\setCJKmainfont{Songti SC}[AutoFakeBold=2]

\definecolor{accent}{HTML}{173F5F}
\definecolor{muted}{HTML}{687782}
\definecolor{line}{HTML}{CBD6DE}

\hypersetup{
  colorlinks=true,
  urlcolor=accent,
  pdfauthor={Cheng Chen},
  pdftitle={陈成其他个人经历与特长}
}

\pagestyle{empty}
\setlength{\parindent}{0pt}
\setlength{\parskip}{0pt}
\setlength{\tabcolsep}{0pt}
\linespread{1.2}
\setlist[itemize]{
  leftmargin=1.2em,
  label=\textcolor{accent}{\small\textbullet},
  itemsep=2pt,
  topsep=2pt,
  parsep=0pt
}

\newcommand{\profilesection}[1]{%
  \vspace{8pt}%
  {\large\bfseries\color{accent}#1}\par
  \vspace{-3pt}\textcolor{line}{\rule{\textwidth}{0.65pt}}\vspace{5pt}%
}

\newcommand{\strength}[2]{%
  \begin{tabularx}{\textwidth}{@{}>{\bfseries\color{accent}}p{3.1cm}X@{}}
    #1 & #2
  \end{tabularx}\vspace{5pt}%
}

\begin{document}

\begin{center}
  {\fontsize{24}{28}\selectfont\bfseries\color{accent}其他个人经历与特长}\\[5pt]
  {\large 陈成}\\[2pt]
  {\color{muted}中国科学院自动化研究所 · 博士生（拟入学）}
\end{center}

\profilesection{团队组织与竞赛经历}

\begin{itemize}
  \item 2026 年中国机器人及人工智能大赛全国一等奖，担任队长。
  \item 2025 年全球校园人工智能算法精英大赛全国二等奖，担任队长。
  \item 2025 年全国大学生嵌入式芯片与系统设计竞赛全国三等奖，担任队长。
  \item 2025 年全国大学生数学建模竞赛、全国大学生统计建模大赛省级三等奖，均担任队长。
  \item 2024 年、2025 年全国大学生数学竞赛省级二等奖。
\end{itemize}

连续担任队长使我形成了清晰的协作方式：先拆分问题和交付物，再根据成员能力分配任务；中期统一代码、数据和文档接口；提交前集中检查结果、表述与材料完整性。这套方法同样适用于科研合作和论文推进。

\profilesection{科研协作与指导经历}

\strength{跨方向适应}{先后参与神经算子可靠性、医学影像蒸馏、开放词汇音视频事件定位和视觉问答研究，能够快速补充任务背景并定位可验证的问题。}
\strength{合作研究}{在 OrthKD 与 OV-OrthKD 中担任共同第一作者，参与方法设计、实验分析和论文写作，重视实验协议与表述的一致性。}
\strength{低年级指导}{在 DeepYOLO 项目中提出核心思路，负责研究设定与方法设计，并指导低年级同学完成实现和实验；论文获国际神经网络联合大会（IJCNN）2026 录用，本人担任通讯作者。}

\profilesection{工作方式}

\strength{研究复盘}{按研究问题组织消融实验，保留配置、日志和结果对应关系；发现结果偏差时，先核对数据划分和运行条件，再定位模型实现。}
\strength{协作接口}{在团队任务中统一代码、数据和文档接口，约定阶段交付物与检查节点，减少不同成员结果无法直接合并的问题。}
\strength{材料把关}{提交前逐项核对术语、数字、图表、实验条件和作者分工，使论文、汇报与附件中的信息保持一致。}

\profilesection{个人特长}

\begin{itemize}
  \item \textbf{问题拆解：}在神经算子研究和 DeepYOLO 合作中，先确定可验证的问题，再拆分方法实现与实验检查项。
  \item \textbf{自主学习：}研究任务跨越科学计算、医学影像和多模态理解，我会围绕具体问题补充数学、优化、信号处理与机器学习知识。
  \item \textbf{沟通协调：}作为多项团队竞赛队长，我会明确接口、时间节点和责任边界，集中记录问题并跟进解决。
  \item \textbf{持续投入：}面对训练异常、复现偏差和写作细节，我会保留排查记录，直到结果具有稳定解释。
\end{itemize}

\vfill
\begin{flushright}
  陈成
\end{flushright}

\end{document}
